\newpage
\section{Vollständige und rohe Ergebnisse des Usability Test}

\begin{figure}[H]
    \centering
    \includegraphics[width=1\linewidth]{}
    \caption{\textcolor{red}{TODO: snippet von der Excel-Tabelle + Formel}}
    \label{fig:Usability_Test:rohe_Daten}
\end{figure}

% *********************************************************************************

\begin{table}[H]
\justifying
\begin{tabular}{|l|c|c|c|c|}
\hline
Aufgabe 1\textcolor{white}{.0}   & $\Bar{x}$ & $\Tilde{x}$ & $\sigma$ & f \\ \hline
alt                              & 16        &  13         &  11      &  0\%   \\ \hline
neu                              & 9,5       &  9          &  3,6     &  9,09\%   \\ \hline
%Differenz   &           &             & -     &  xxx   \\ \hline
%\end{tabular}
%\begin{tabular}{|l|c|c|c|c|}
\hline
Aufgabe 2\textcolor{white}{.0}   & $\Bar{x}$ & $\Tilde{x}$ & $\sigma$ & f \\ \hline
alt                              & 14        &  10         &  9,1  &  0\%   \\ \hline
neu                              & 6,5       &  6,5        &  2,3  &  30\%   \\ \hline
%Differenz   &           &             & -     &  xxx   \\ \hline
%\end{tabular}
%\begin{tabular}{|l|c|c|c|c|}
\hline
Aufgabe 3.1   & $\Bar{x}$ & $\Tilde{x}$ & $\sigma$ & f \\ \hline
alt           & 16        &  15         &  6,2  &  0\%   \\ \hline
neu           & 12        &  11         &  6,4  &  36,36\%   \\ \hline
%Differenz   &           &             & -     &  xxx   \\ \hline
\end{tabular}
\begin{tabular}{|l|c|c|c|c|}
\hline
Aufgabe 3.2   & $\Bar{x}$ & $\Tilde{x}$ & $\sigma$ & f \\ \hline
alt           & 9         &  8          &  5,4  &  0\%   \\ \hline
neu           & 5,4       &  4          &  3    &  36,36\%   \\ \hline
%Differenz   &           &             & -     &  xxx   \\ \hline
%\end{tabular}
%\begin{tabular}{|l|c|c|c|c|}
\hline
Aufgabe 4.1   & $\Bar{x}$ & $\Tilde{x}$ & $\sigma$ & f \\ \hline
alt           & 21        &  18         &  7,4  &  45,45\%   \\ \hline
neu           & 15        &  12         &  8,8  &  9,09\%   \\ \hline
%Differenz   &           &             & -     &  xxx   \\ \hline
%\end{tabular}
%\begin{tabular}{|l|c|c|c|c|}
\hline
Aufgabe 4.2   & $\Bar{x}$ & $\Tilde{x}$ & $\sigma$ & f \\ \hline
alt           & 12        &  11         &  7,7  &  18,18\%   \\ \hline
neu           & 8,5       &  7          &  4,4  &  9,09\%   \\ \hline
%Differenz   &           &             & -     &  xxx   \\ \hline
\end{tabular}

\begin{tabular}{ll}
\multicolumn{2}{l}{verwendete Formelzeichen:} \\ \hline
$\Bar{x}$          & Mittelwert               \\
$\Tilde{x}$        & Median                   \\
$\sigma$           & Standardabweichung       \\
f                  & Fehlerquote
\end{tabular}
\caption{Zusammenfassung der Ergebnisse des Usability-Tests (Weiterverarbeitung der Daten aus \autoref{fig:Usability_Test:rohe_Daten})}
\label{tab:usability_test:mittelwerte}
\end{table}

% *********************************************************************************

\begin{table}[H]
\begin{tabular}{|l|c|c|c|c|c|c|}
\hline
                                          & Frage 1 & Frage 2 & Frage 3.1 & Frage 3.2 & Frage 4.1 & Frage 4.2 \\ \hline
$\varnothing t_{neu} - \varnothing t_{alt}$ & -6,55   & -7,14   & -4,46     & -3,64     & -6,18     & -3,81     \\ \hline
T-Test                                    & 0,0231  & 0,0186  & 0,0579    & 0,0506    & 0,0554    & 0,0826    \\ \hline
\end{tabular}
\caption{Usability Test - Differenz der Zeitmittelwerte und T-Test-Ergebnisse (Weiterverarbeitung der Daten aus \autoref{tab:usability_test:mittelwerte} und \autoref{fig:Usability_Test:rohe_Daten})}
\label{tab:usability_test:Differenz+ttest}
\end{table}

\textcolor{red}{TODO: weitere T-Tests von Aufgaben 2+3}\\
?? hier? oder war ddas bei der Nutzerbefragung?